% ============================================================
%  第一章  概率、图与因果模型导论  （第 1 节）
% ============================================================
\chapter{概率、图与因果模型导论}

\begin{flushright}
\kaishu\small
偶然生出思绪，偶然又将其抹去。\\[2pt]
——帕斯卡（Pascal，1670）
\end{flushright}

\section{概率论导引}

\subsection{为何需要概率？}

因果性（causality）意味着合乎规律的必然，概率则意味着例外、疑虑与不规则。尽管如此，从概率分析入手、甚至着重于概率分析来讨论因果性，仍有两个强有力的理由：一个相当直白，另一个则较为微妙。

简单的理由来自一个观察：因果话语常常出现在充满不确定性的情境之中。我们会说“鲁莽驾驶导致车祸”，或者“你之所以挂科是因为懒惰”（Suppes 1970），同时心里很清楚，这些前件只是使后果更可能发生，而非使之必然发生。任何试图容纳此类话语的因果理论，都必须用一种能够区分各种可能性程度的语言来表述——也就是概率的语言。与此相关，我们还注意到，概率论如今已是大多数使用因果建模的学科的官方数学语言，包括经济学、流行病学、社会学和心理学。在这些学科中，研究者关心的不只是因果联系的有无，还包括联系的相对强度，以及如何从带有噪声的观测中推断出这些联系。概率论辅以统计分析方法，既提供了应对这类观测的原则，也提供了从中作出推断的手段。

更微妙的理由则涉及这样一个事实：即便是自然语言中最斩钉截铁的因果表述，也允许例外；而这些例外若交由标准的确定性逻辑规则处理，就可能造成严重的困难。试看下面两个看似成立的前提：

\begin{enumerate}[label=\arabic*.]
\item 只要我家屋顶湿了，邻居家的屋顶就湿了。
\item 如果我用水管冲我家的屋顶，它就会湿。
\end{enumerate}

照字面理解，这两个前提会推出一个荒谬的结论：只要我冲洗自家屋顶，邻居家的屋顶就会湿。

这类悖论通常被归因于语言的有限粒度，它体现为前提 1 中隐含的诸多例外。事实上，只要我们不怕麻烦地把这些例外显式说出来，悖论便消失了，例如写成：

\begin{enumerate}[label=\arabic*$^{*}$.]
\item 只要我家屋顶湿了，邻居家的屋顶就湿了——除非他家屋顶盖着塑料布，或者我正用水管冲我家屋顶，等等。
\end{enumerate}

概率论恰恰擅长容忍未经言明的例外，因而使我们能够聚焦于因果性的主要问题，而不必纠缠于此类悖论。

后文将会看到，容忍例外只解决了与因果性相关的部分问题。其余问题——包括推断、干预、识别、推演、混杂、反事实与解释——才是本书的主题。用概率的语言来刻画这些问题，是为了强调它们跨语言的普遍性。第七章将把这些问题的表述改换为确定性逻辑的语言，并把概率仅仅当作表达“对未观测事实之不确定”的一种方式。

\subsection{概率论的基本概念}

本书的讨论大多集中于具有有限个离散变量的系统，因而只需概率论中最基本的记号与初等概念。向连续变量的推广将予以概述，但不作完全一般性的展开。希望补充数学工具的读者，可以研读该领域众多优秀的教科书，例如 Feller (1950)、Hoel 等 (1971)，或 Suppes (1970) 的附录。本节简要概述初等概率概念，主要依据 Pearl (1988b)，并特别着重于贝叶斯推断及其与人类不确定性推理心理学的关联——这一侧重在标准教科书中通常是缺失的。

我们遵循概率的贝叶斯解释：概率刻画的是对世界中之事件的信念度，数据则被用来加强、更新或削弱这些信念度。在这一形式体系中，信念度被赋予某种语言中的命题（取真值或假值的语句），并按概率演算的规则加以组合与运算。我们不区分语句性命题与这些命题所表示的实际事件。例如，若 $A$ 表示语句“泰德·肯尼迪将在 2012 年争取总统提名”，则 $P(A\mid K)$ 表示某人在知识体 $K$ 之下对 $A$ 所描述事件的主观信念，而 $K$ 可能包含此人关于美国政治的假设、肯尼迪的某些公开声明，以及对其年龄与个性的评估。在书写概率表达式时，我们常常只写 $P(A)$ 而略去符号 $K$。然而，当背景信息发生变化时，我们就需要明确指出支撑信念的假设，并把 $K$（或其部分要素）显式写出来。

在贝叶斯形式体系中，信念测度服从概率演算的三条基本公理：
\begin{gather}
0 \le P(A) \le 1, \tag{1.1}\\
P(\text{必然命题}) = 1, \tag{1.2}\\
P(A \text{ 或 } B) = P(A) + P(B)\quad\text{若 } A \text{ 与 } B \text{ 互斥}. \tag{1.3}
\end{gather}
第三条公理说的是：赋予任一事件集合的信念，等于赋予其各不相交组成部分的信念之和。由于任一事件 $A$ 都可写成联合事件 $(A\wedge B)$ 与 $(A\wedge\neg B)$ 的并，故其概率由下式给出\footnote{符号 $\wedge,\ \vee,\ \neg,\ \Rightarrow$ 分别表示逻辑联结词“与”“或”“非”“蕴含”。}：
\begin{equation}
P(A) = P(A,B) + P(A,\neg B), \tag{1.4}
\end{equation}
其中 $P(A,B)$ 是 $P(A\wedge B)$ 的简写。更一般地，若 $B_i,\ i=1,2,\dots,n$ 是一组穷尽且互斥的命题（称为一个\emph{划分}，或一个\emph{变量}），则 $P(A)$ 可由 $P(A,B_i),\ i=1,2,\dots,n$ 求和得到：
\begin{equation}
P(A) = \sum_i P(A,B_i), \tag{1.5}
\end{equation}
此式即所谓“全概率公式”。对一切 $B_i$ 求和这一运算也称为“对 $B$ 边缘化”，所得概率 $P(A)$ 称为 $A$ 的\emph{边缘概率}。例如，事件 $A$“两枚骰子的点数相同”的概率，可通过对联合事件 $(A\wedge B_i),\ i=1,2,\dots,6$ 求和得到，其中 $B_i$ 表示命题“第一枚骰子的点数为 $i$”。于是
\begin{equation}
P(A) = \sum_i P(A,B_i) = 6\times\frac{1}{36} = \frac{1}{6}. \tag{1.6}
\end{equation}

由 \textup{(1.2)} 与 \textup{(1.4)} 可直接得到：一个命题与其否定所分得的信念总和必为 1，
\begin{equation}
P(A) + P(\neg A) = 1, \tag{1.7}
\end{equation}
因为这两个陈述中必有一个为真。

贝叶斯形式体系中的基本表达式是关于\emph{条件概率}的陈述，例如 $P(A\mid B)$，它给出在假设 $B$ 以绝对确定性已知的条件下对 $A$ 的信念。若 $P(A\mid B)=P(A)$，我们就说 $A$ 与 $B$ \emph{独立}，因为得知 $B$ 为真并未改变我们对 $A$ 的信念。若 $P(A\mid B,C)=P(A\mid C)$，我们就说在给定 $C$ 的条件下 $A$ 与 $B$ \emph{条件独立}；也就是说，一旦知道了 $C$，再得知 $B$ 也不会改变我们对 $A$ 的信念。

与以联合事件定义条件概率的传统做法相反，
\begin{equation}
P(A\mid B) = \frac{P(A,B)}{P(B)}, \tag{1.8}
\end{equation}
贝叶斯派哲学家认为条件关系比联合事件更为基本——也就是说，它更契合人类知识的组织方式。依此观点，$B$ 充当指向某个知识语境或框架的指针，而 $A\mid B$ 表示在 $B$ 所指定的语境中的事件 $A$（例如，在疾病 $B$ 语境中的症状 $A$）。因此，经验知识总是以条件概率陈述的形式被编码，而对联合事件的信念（如果确有需要）则由这些陈述通过乘积计算出来：
\begin{equation}
P(A,B) = P(A\mid B)P(B), \tag{1.9}
\end{equation}
它与 \textup{(1.8)} 等价。例如，在 \textup{(1.6)} 中直接评估 $P(A,B_i)$ 多少有些别扭。这种评估背后的心理过程预设了两个结果是独立的；为使该假设显式化，联合事件（点数相同，$B_i$）的概率应当由条件事件（点数相同 $\mid B_i$）经乘积得到：
\[
P(\text{点数相同}\mid B_i)\,P(B_i) = P(\text{第二枚骰子点数为 } i\mid B_i)\,P(B_i) = \frac{1}{6}\times\frac{1}{6} = \frac{1}{36}.
\]

如同 \textup{(1.5)}，任一事件 $A$ 的概率都可以先对一组穷尽且互斥的事件 $B_i,\ i=1,2,\dots,n$ 取条件，再求和得到：
\begin{equation}
P(A) = \sum_i P(A\mid B_i)P(B_i). \tag{1.10}
\end{equation}
这一分解为假设性推理、即“基于假设的”推理提供了基础。它表明，对任一事件 $A$ 的信念，是对 $A$ 得以实现的各个不同途径之信念的加权和。例如，若要计算第一枚骰子的点数 $X$ 大于第二枚骰子点数 $Y$ 的概率，我们可以把事件 $A: X>Y$ 对 $X$ 的一切可能取值取条件，得到
\begin{equation}
\begin{aligned}
P(A) &= \sum_{i=1}^{6} P(Y<X \mid X=i)\,P(X=i)
= \sum_{i=1}^{6} P(Y<i)\,\frac{1}{6}
= \sum_{i=1}^{6}\sum_{j=1}^{i-1} P(Y=j)\,\frac{1}{6}\\
&= \frac{1}{6}\sum_{i=2}^{6}\frac{i-1}{6} = \frac{5}{12}.
\end{aligned} \tag{1.11}
\end{equation}
值得再次强调：像 \textup{(1.10)} 这样的公式，始终是在某个更大的语境 $K$ 中理解的，$K$ 定义了被视为共同知识的那些假设（例如掷骰子的公平性）。式 \textup{(1.10)} 实际上是下述陈述的简写：
\begin{equation}
P(A\mid K) = \sum_i P(A\mid B_i,K)P(B_i\mid K). \tag{1.12}
\end{equation}
此式成立，是因为每个条件概率 $P(A\mid K)$ 本身就是一个真正的概率函数，因而满足 \textup{(1.10)}。

乘积规则（式 \textup{(1.9)}）的另一有用推广是\emph{链式法则}。它说的是：若有一组 $n$ 个事件 $E_1,E_2,\dots,E_n$，则联合事件 $(E_1,E_2,\dots,E_n)$ 的概率可写成 $n$ 个条件概率之积：
\begin{equation}
P(E_1,E_2,\dots,E_n) = P(E_n \mid E_{n-1},\dots,E_2,E_1)\cdots P(E_2\mid E_1)P(E_1). \tag{1.13}
\end{equation}
该乘积可按任意方便的次序反复应用 \textup{(1.9)} 得到。

贝叶斯推断的核心在于著名的反演公式
\begin{equation}
P(H\mid e) = \frac{P(e\mid H)P(H)}{P(e)}, \tag{1.14}
\end{equation}
它表明：在获得证据 $e$ 之后赋予假设 $H$ 的信念，等于我们原有的信念 $P(H)$ 乘以“若 $H$ 为真则 $e$ 将出现”的似然 $P(e\mid H)$。$P(H\mid e)$ 有时称为\emph{后验概率}（或简称后验），$P(H)$ 称为\emph{先验概率}（或简称先验）。式 \textup{(1.14)} 的分母 $P(e)$ 几乎不必考虑，因为它只是一个归一化常数
\[
P(e) = P(e\mid H)P(H) + P(e\mid\neg H)P(\neg H),
\]
可由“$P(H\mid e)$ 与 $P(\neg H\mid e)$ 之和为 1”这一要求算出。

从形式上看，\textup{(1.14)} 也许会被斥为条件概率定义的同义反复
\[
P(A\mid B) = \frac{P(A,B)}{P(B)} \quad\text{与}\quad P(B\mid A) = \frac{P(A,B)}{P(A)},
\]
但贝叶斯主观主义者把 \textup{(1.14)} 视为“依据证据更新信念”的规范性规则。换言之，条件概率固然可以被看作纯粹的数学构造，贝叶斯派却视之为语言的原初概念，视为对英语表达“……，已知 $A$”的忠实翻译。据此，上述定义式就不是一个定义，而是英语表达之间可由经验检验的关系。它断言：一个人发现 $A$ 之后赋予 $B$ 的信念，绝不会低于发现 $A$ 之前赋予 $A\wedge B$ 的信念；并且这两个信念之比，将随“发现 $A$ 带来的惊讶程度”$[P(A)]^{-1}$ 成比例增大。

式 \textup{(1.14)} 的重要性在于，它把人们往往难以评估的量 $P(H\mid e)$，用通常可以直接从经验知识中取得的量表达出来。例如，隔壁赌桌上有人报出结果“十二点”，我们想知道他掷的是双骰还是转的轮盘；关于赌具的模型立刻给出 $P(\text{十二}\mid\text{骰子})$ 与 $P(\text{十二}\mid\text{轮盘})$：前者为 $1/36$，后者为 $1/38$。同样，我们可以通过估计赌场中轮盘与骰桌的数量来判断先验 $P(\text{骰子})$ 与 $P(\text{轮盘})$。而直接判断 $P(\text{骰子}\mid\text{十二})$ 则要困难得多；只有在这家赌场受过训练的行家才能可靠地作出这种判断。

为完成这一简要导引，我们还须讨论\emph{概率模型}（又称概率空间）的概念。概率模型是信息的编码方式，它使我们能够依据 \textup{(1.1)}—\textup{(1.3)} 的公理计算每一个合式语句 $S$ 的概率。从一组原子命题 $A,B,C,\dots$ 出发，合式语句的集合由涉及这些命题的一切布尔公式构成，例如
\[
S = (A\wedge B)\vee\neg C .
\]
传统的概率模型指定方法采用\emph{联合分布函数}：它给语言中的每个基本事件（基本事件是这样一个合取式，其中每个原子命题或其否定恰好出现一次）赋予非负权重，并使权重之和为 1。例如，若有三个原子命题 $A,B,C$，则联合分布函数应给全部八种组合
\[
(A\wedge B\wedge C),\ (A\wedge B\wedge\neg C),\ \dots,\ (\neg A\wedge\neg B\wedge\neg C)
\]
赋予非负权重，并使八个权重之和为 1。

读者或许已经认出，基本事件的集合就是概率教科书中的样本空间。例如，若 $A,B,C$ 分别对应“第 1、2、3 枚硬币正面朝上”这三个命题，则样本空间由集合 $\{HHH,HHT,HTH,\dots,TTT\}$ 构成。事实上，把对应于基本事件的合取式视为点（或世界、或配置），而把其他公式视为由这些点组成的集合，有时是很方便的。由于每个布尔公式都可表示为基本事件的析取，而基本事件互斥，我们总能用可加性公理（式 \textup{(1.3)}）计算 $P(S)$。条件概率同样可用 \textup{(1.8)} 计算。因此，任何联合概率函数都表示一个完整的概率模型。

联合分布函数是极为重要的数学构造。它们使我们能够迅速判断：现有信息是否足以确定一个完整的概率模型，已有信息是否自洽，以及在何处还需要补充信息。判据很简单：(i) 现有信息是否足以唯一确定论域中每个基本事件的概率；(ii) 这些概率之和是否为 1。

然而在实践中，联合分布函数很少被显式给出。在连续随机变量的分析中，分布函数由代数表达式给出，例如描述正态分布或指数分布的表达式；对于离散变量，则发展出了间接的表示方法，即由小变量组之间的局部关系推断整体分布。图模型是这类表示中最流行的一种，也是本书通篇讨论的基础。其用法与形式刻画将在随后几节中讨论。

\subsection{预测性支持与诊断性支持的结合}

贝叶斯法则（式 \textup{(1.14)}）的精髓，可以方便地用\emph{几率}与\emph{似然比}这两个参数来刻画。将 \textup{(1.14)} 除以其关于 $\neg H$ 的互补形式 $P(\neg H\mid e)$，得到
\begin{equation}
\frac{P(H\mid e)}{P(\neg H\mid e)} = \frac{P(e\mid H)}{P(e\mid\neg H)}\,\frac{P(H)}{P(\neg H)}. \tag{1.15}
\end{equation}
定义 $H$ 的\emph{先验几率}为
\begin{equation}
O(H) = \frac{P(H)}{P(\neg H)} = \frac{P(H)}{1-P(H)}, \tag{1.16}
\end{equation}
\emph{似然比}为
\begin{equation}
L(e\mid H) = \frac{P(e\mid H)}{P(e\mid\neg H)}, \tag{1.17}
\end{equation}
则\emph{后验几率}
\begin{equation}
O(H\mid e) = \frac{P(H\mid e)}{P(\neg H\mid e)} \tag{1.18}
\end{equation}
由下式之积给出：
\begin{equation}
O(H\mid e) = L(e\mid H)\,O(H). \tag{1.19}
\end{equation}
于是，贝叶斯法则告诉我们：基于原有知识 $K$ 与观测证据 $e$ 而对假设 $H$ 持有的总体信念强度，应当是两个因子之积——先验几率 $O(H)$ 与似然比 $L(e\mid H)$。\textbf{第一个因子度量背景知识单独给予 $H$ 的预测性支持（或前瞻性支持），第二个因子则代表实际观测到的证据给予 $H$ 的诊断性支持（或回溯性支持）。}\footnote{在流行病学中，若 $H$ 表示暴露、$e$ 表示疾病，则似然比 $L$ 被称为“风险比”（Rothman and Greenland 1998, p.~50）。此时式 \textup{(1.18)} 给出的是“患疾病 $e$ 者曾暴露于 $H$”的几率。}

严格说来，似然比 $L(e\mid H)$ 可能依赖于默会知识库 $K$ 的内容。然而，贝叶斯技术的力量主要来自这样一个事实：在因果推理中，关系 $P(e\mid H)$ 是相当局部的——给定 $H$ 为真，$e$ 的概率可以自然地估计出来，因为它通常并不依赖于知识库中许多其他命题。例如，一旦确定某病人患有疾病 $H$，估计她将出现某症状 $e$ 的概率便是自然的。医学知识的组织建立在这样一个范式之上：症状是疾病的稳定特征，因而应当相当独立于其他因素，例如流行状况、既往病史和诊断设备的故障。正因如此，$P(e\mid H)$ 而非 $P(H\mid e)$ 才是贝叶斯分析中的原子关系。前者具有类似于逻辑规则的模块性：它们表达了对“若 $H$ 则 $e$”这类规则的置信程度，而这种置信不因知识库中还存在哪些其他规则或事实而改变。

\begin{example}[入室盗窃与警报]
设想某夜你被刺耳的防盗警报惊醒。你相信确实发生了入室盗窃的信念度是多少？为便于说明，我们作如下判断：(a) 入室盗窃触发警报系统的机会为 95\%，即 $P(\text{警报}\mid\text{盗窃})=0.95$；(b) 根据以往误报情况，警报由盗窃以外的机制触发的机会很小，为 1\%，即 $P(\text{警报}\mid\text{无盗窃})=0.01$；(c) 以往犯罪统计表明，某住宅在某夜被窃的机会为万分之一，即 $P(\text{盗窃})=10^{-4}$。

把这些假设用 \textup{(1.19)} 结合起来，得到
\[
O(\text{盗窃}\mid\text{警报}) = L(\text{警报}\mid\text{盗窃})\,O(\text{盗窃}) = \frac{0.95}{0.01}\times\frac{10^{-4}}{1-10^{-4}} = 0.0095 .
\]
于是，由
\begin{equation}
P(A) = \frac{O(A)}{1+O(A)} \tag{1.20}
\end{equation}
可得
\[
P(\text{盗窃}\mid\text{警报}) = \frac{0.0095}{1+0.0095} = 0.00941 .
\]
可见，警报证据给予盗窃假设的回溯性支持，把它的信念度提高了近一百倍——从万分之一提高到万分之九十四点一。盗窃信念仍低于 1\%，这并不奇怪，因为该系统差不多每三个月就会误报一次。注意，我们并不需要估计 $P(\text{警报}\mid\text{盗窃})$ 与 $P(\text{警报}\mid\text{无盗窃})$ 的绝对值；计算中只用到二者之比，因此也可以直接估计这个比值。
\end{example}

\subsection{随机变量与期望}

所谓\emph{变量}，我们指某种属性、测量或询问，它可以从一个指定论域中取若干可能结果（即取值）之一。若我们对一个变量可能取得的各个值附有信念（即概率），就称该变量为\emph{随机变量}。\footnote{这是对教科书定义的一点小推广。教科书把随机变量定义为从样本空间（例如基本事件集）到实数轴的映射；在我们的定义中，映射的靶是任意一个称为“取值”的对象集合，其中元素未必有序。}例如，我明天所穿鞋子的颜色是一个名为“颜色”的随机变量，其取值来自论域 $\{\text{黄},\text{绿},\text{红},\dots\}$。

我们的分析大多涉及一个有限的随机变量集合 $V$（随机变量也称划分），其中每个变量 $X\in V$ 取值于有限论域 $D_X$。我们用大写字母（如 $X,Y,Z$）表示变量名，用小写字母（$x,y,z$）表示相应变量所取的某个具体值。例如，若 $X$ 表示物体的颜色，则 $x$ 指从集合 $\{\text{黄},\text{绿},\text{红},\dots\}$ 中任取的一个元素。显然，命题 $X=\text{黄}$ 描述了一个事件，即满足“该物体颜色为黄”这一命题的那些可能事态的集合。同样，每个变量 $X$ 都可视为对世界状态的一种划分，因为语句 $X=x$ 定义了若干穷尽且互斥的状态集合，$x$ 的每个取值对应其中一个。

在本书的大部分讨论中，我们不在记号上区分变量与变量集合，因为一个变量集本质上定义了一个复合变量，其论域是集合中各成员论域的笛卡儿积。因此，若 $Z$ 表示集合 $\{X,Y\}$，则 $z$ 表示数对 $(x,y)$，其中 $x\in D_X$、$y\in D_Y$。当需要特别强调变量与变量集之别时，我们用带下标的字母（如 $X_1,X_2,\dots,X_n$ 或 $V_1,V_2,\dots,V_n$）表示各个变量。

我们将一贯地用简写 $P(x)$ 表示概率 $P(X=x)$，$x\in D_X$。同样，若 $Z$ 表示 $\{X,Y\}$，则 $P(z)$ 定义为
\[
P(z) \triangleq P(Z=z) = P(X=x,\,Y=y),\qquad x\in D_X,\ y\in D_Y .
\]
当随机变量 $X$ 取实数值时，称 $X$ 为\emph{实随机变量}；此时可定义 $X$ 的均值即期望为
\begin{equation}
E(X) \triangleq \sum_x x\,P(x), \tag{1.21}
\end{equation}
并定义在给定事件 $Y=y$ 之下 $X$ 的条件均值为
\begin{equation}
E(X\mid y) \triangleq \sum_x x\,P(x\mid y). \tag{1.22}
\end{equation}
$X$ 的任一函数 $g$ 的期望定义为
\begin{equation}
E[g(X)] \triangleq \sum_x g(x)P(x). \tag{1.23}
\end{equation}
其中，函数 $g(X)=(X-E(X))^2$ 格外受重视，其期望称为 $X$ 的\emph{方差}，记作
\[
\sigma_X^2 \triangleq E[(X-E(X))^2].
\]
在给定观测 $Y=y$ 的条件下，条件均值 $E(X\mid Y=y)$ 是 $X$ 的最佳估计，其含义是它使期望平方误差
\[
\sum_{x} (x-x')^2 P(x\mid y)
\]
对一切可能的估计 $x'$ 取最小。

二元函数 $g(X,Y)$ 的期望需要联合概率 $P(x,y)$，其定义为
\[
E[g(X,Y)] \triangleq \sum_{x,y} g(x,y)P(x,y)
\]
（参看式 \textup{(1.23)}）。特别重要的是乘积的期望，即 $g(X,Y)=(X-E(X))(Y-E(Y))$，它称为 $X$ 与 $Y$ 的\emph{协方差}
\[
\sigma_{XY} \triangleq E[(X-E(X))(Y-E(Y))],
\]
协方差常被归一化，得到\emph{相关系数}
\[
r_{XY} = \frac{\sigma_{XY}}{\sigma_X \sigma_Y}
\]
以及（$X$ 对 $Y$ 的）\emph{回归系数}
\[
r_{YX} \triangleq r_{XY}\frac{\sigma_X}{\sigma_Y} = \frac{\sigma_{XY}}{\sigma_Y^2}.
\]
给定 $Z=z$ 的条件方差、条件协方差与条件相关系数，可用条件分布 $P(x,y\mid z)$ 求期望，以类似方式定义。特别地，给定 $Z=z$ 的条件相关系数定义为
\begin{equation}
r_{XY\mid z} = \frac{\sigma_{XY\mid z}}{\sigma_{X\mid z}\,\sigma_{Y\mid z}}. \tag{1.24}
\end{equation}
另有一些专属于正态分布的性质，将在第五章（\S 5.2.1）回顾。

上述定义适用于\emph{离散}随机变量，即在实数轴上取有限或可数无限个值的变量。期望与相关的讨论更常用于\emph{连续}随机变量，后者由密度函数 $f(x)$ 刻画，其定义为：对任意满足 $a<b$ 的实数 $a,b$，
\[
P(a<X<b) = \int_a^b f(x)\,dx .
\]
若 $X$ 为离散型，则只要按如下对应关系理解积分，$f(x)$ 便与概率函数 $P(x)$ 一致：
\begin{equation}
\int_{-\infty}^{\infty} f(x)\,dx \;\longleftrightarrow\; \sum_x P(x). \tag{1.25}
\end{equation}
习惯于连续分析的读者，在本书中遇到求和号时应记住这一对应。例如，连续随机变量 $X$ 的期望可由 \textup{(1.21)} 改写为
\[
E(X) = \int_{-\infty}^{\infty} x f(x)\,dx,
\]
方差、相关等亦作类似改写。

下面我们定义变量间的条件独立关系——这是因果建模中的一个核心概念。

\subsection{条件独立与图拟}

\begin{definition}[条件独立]
设 $V=\{V_1,V_2,\dots\}$ 是一个有限变量集，$P(\cdot)$ 是 $V$ 中变量上的联合概率函数，$X,Y,Z$ 表示 $V$ 中任意三个变量子集。若
\begin{equation}
P(x\mid y,z) = P(x\mid z)\qquad\text{只要 } P(y,z)>0, \tag{1.26}
\end{equation}
则称集合 $X$ 与 $Y$ 在给定 $Z$ 的条件下\emph{条件独立}。
\end{definition}

用文字说，一旦知道了 $Z$，再得知 $Y$ 的取值并不会为 $X$ 提供额外信息。（打个比方，$Z$ 把 $X$ 与 $Y$ “屏蔽”开了。）

式 \textup{(1.26)} 是下述说法的简略表达：对集合 $X$ 中变量的任一配置 $x$，以及对 $Y$ 与 $Z$ 中变量的任一组满足 $P(Y=y,Z=z)>0$ 的配置 $y,z$，都有
\begin{equation}
P(X=x\mid Y=y,Z=z) = P(X=x\mid Z=z). \tag{1.27}
\end{equation}
我们将采用 Dawid (1979) 的记号 $(X\indep Y\mid Z)_P$，或简写为 $(X\indep Y\mid Z)$，表示给定 $Z$ 时 $X$ 与 $Y$ 条件独立；于是
\begin{equation}
(X\indep Y\mid Z)_P \quad\text{当且仅当}\quad P(x\mid y,z)=P(x\mid z), \tag{1.28}
\end{equation}
对一切满足 $P(y,z)>0$ 的取值 $x,y,z$ 成立。无条件独立（也称边缘独立）记作 $(X\indep Y\mid\varnothing)$，即
\begin{equation}
(X\indep Y\mid\varnothing)_P \quad\text{当且仅当}\quad P(x\mid y)=P(x)\ \text{只要}\ P(y)>0. \tag{1.29}
\end{equation}
注意，$(X\indep Y\mid Z)$ 蕴含 $X$ 与 $Y$ 中所有变量对之间条件独立，即对一切 $V_i\in X$、$V_j\in Y$ 有 $(V_i\indep V_j\mid Z)$，但反之未必成立。

下面列出条件独立关系 $(X\indep Y\mid Z)$ 所满足的性质（不完全清单）：
\begin{align*}
&\text{对称性：} && (X\indep Y\mid Z)\ \Rightarrow\ (Y\indep X\mid Z);\\
&\text{分解性：} && (X\indep YW\mid Z)\ \Rightarrow\ (X\indep Y\mid Z);\\
&\text{弱联合：} && (X\indep YW\mid Z)\ \Rightarrow\ (X\indep Y\mid ZW);\\
&\text{收缩性：} && (X\indep Y\mid Z)\ \wedge\ (X\indep W\mid ZY)\ \Rightarrow\ (X\indep YW\mid Z);\\
&\text{交性质：} && (X\indep W\mid ZY)\ \wedge\ (X\indep Y\mid ZW)\ \Rightarrow\ (X\indep YW\mid Z).
\end{align*}
（交性质在严格正的概率分布中成立。）

这些性质的证明可由 \textup{(1.28)} 及概率论基本公理初等地推出。\footnote{这些性质最先由 Dawid (1979) 与 Spohn (1980) 以略有不同的形式引入，并由 Pearl and Paz (1987) 独立提出，用以刻画图与信息相关性之间的关系。Geiger and Pearl (1993) 给出了深入分析。}它们被 Pearl and Paz (1987) 与 Geiger 等 (1990) 称为\emph{图拟公理}（graphoid axioms），并被证明在多种解释下支配着信息相关性的概念（Pearl 1988b）。例如，在图论中，若把 $(X\indep Y\mid Z)$ 解释为“从节点子集 $X$ 到节点子集 $Y$ 的所有路径都被节点子集 $Z$ 截断”，则这些性质成立。

图拟公理的直观解释如下（Pearl 1988b, p.~85）。对称性公理说的是：在任一知识状态 $Z$ 下，若 $Y$ 不能告诉我们关于 $X$ 的新信息，则 $X$ 也不能告诉我们关于 $Y$ 的新信息。分解性公理断言：若两条合并的信息被认为与 $X$ 无关，则其中每一条单独的信息也与 $X$ 无关。弱联合公理说的是：获知无关信息 $W$，不会使原本无关的信息 $Y$ 变得与 $X$ 相关。收缩性公理说的是：若在获知某条无关信息 $Y$ 之后我们判定 $W$ 与 $X$ 无关，那么 $W$ 在获知 $Y$ 之前就应当是无关的。弱联合与收缩合在一起意味着：无关信息不应改变系统中其他命题的相关性状态——原本相关的仍然相关，原本无关的仍然无关。交性质公理说的是：若在已知 $W$ 时 $Y$ 与 $X$ 无关，且在已知 $Y$ 时 $W$ 与 $X$ 无关，则 $W$ 与 $Y$（以及二者的组合）都与 $X$ 无关。
